\section{Training MoEs：负载均衡、随机路由与系统效率}

本节处理 MoE 训练的核心难题：既要稀疏计算，又要专家负载均衡，还要保持可训练稳定。


\begin{warningbox}{MoE 训练的核心矛盾}
MoE 想要稀疏激活来省 FLOPs，但系统效率要求 experts 被均匀使用。如果 router 总把 token 发给少数 experts，就会出现负载不均、通信拥塞、部分 experts 欠训练。
\end{warningbox}

\begin{figure}[H]
\centering
\includegraphics[width=0.88\textwidth]{slide-035.jpg}
\caption{Slide 36: How do we train MoEs?}
\figsource{CS336 Spring 2026 Lecture 4 官方幻灯片。}
\end{figure}

\noindent\textbf{展开说明：}Major challenge: we need sparsity for training-time efficiency…


\begin{figure}[H]
\centering
\includegraphics[width=0.88\textwidth]{slide-036.jpg}
\caption{Slide 37: RL for MoEs}
\figsource{CS336 Spring 2026 Lecture 4 官方幻灯片。}
\end{figure}

\noindent\textbf{读图说明：}把 expert 选择视为离散 action 后，可以用 REINFORCE 根据下游 loss 学习 routing policy。它理论上允许直接优化非可微选择，但梯度方差高、credit assignment 困难，还要设计 baseline 与探索噪声。课程结论不是“RL 无效”，而是收益不足以稳定抵消复杂度；在已有 top-$k$、辅助负载损失和可微近似表现良好时，RL routing 尚未形成明确默认值。


\begin{figure}[H]
\centering
\includegraphics[width=0.88\textwidth]{slide-037.jpg}
\caption{Slide 38: Stochastic approximations}
\figsource{CS336 Spring 2026 Lecture 4 官方幻灯片。}
\end{figure}

\begin{figure}[H]
\centering
\includegraphics[width=0.88\textwidth]{slide-038.jpg}
\caption{Slide 39：stochastic approximation 的第二步推导与可训练近似。}
\figsource{CS336 Spring 2026 Lecture 4 官方幻灯片。}
\end{figure}

\noindent\textbf{展开说明：}From Shazeer et al 2017 – routing decisions are stochastic with gaussian perturbations.



\noindent\textbf{展开说明：}Stochastic jitter in Fedus et al 2022. This does a uniform multiplicative perturbation for the


\begin{figure}[H]
\centering
\includegraphics[width=0.88\textwidth]{slide-039.jpg}
\caption{Slide 40: Heuristic balancing losses}
\figsource{CS336 Spring 2026 Lecture 4 官方幻灯片。}
\end{figure}

\noindent\textbf{展开说明：}Another key issue – systems efficiency requires that we use experts evenly..

\begin{knowledgebox}{读图：Slide 40 应该怎么看}
这页是机制或证据页。先看它比较的是 compute、参数量、routing、负载均衡还是通信；再看结论支持的是质量提升、训练速度、推理成本还是系统复杂度。MoE 和 attention alternatives 的图常把模型结构和系统代价混在一起，读图时要分清“算法机制”和“硬件执行成本”。
\end{knowledgebox}


\begin{figure}[H]
\centering
\includegraphics[width=0.88\textwidth]{slide-040.jpg}
\caption{Slide 41: Example from deepseek (v1-2)}
\figsource{CS336 Spring 2026 Lecture 4 官方幻灯片。}
\end{figure}

\noindent\textbf{展开说明：}Per-expert balancing – same as the switch transformer

\begin{knowledgebox}{读图：Slide 41 应该怎么看}
这页是机制或证据页。先看它比较的是 compute、参数量、routing、负载均衡还是通信；再看结论支持的是质量提升、训练速度、推理成本还是系统复杂度。MoE 和 attention alternatives 的图常把模型结构和系统代价混在一起，读图时要分清“算法机制”和“硬件执行成本”。
\end{knowledgebox}


\begin{figure}[H]
\centering
\includegraphics[width=0.88\textwidth]{slide-041.jpg}
\caption{Slide 42: DeepSeek v3 variation – per-expert biases}
\figsource{CS336 Spring 2026 Lecture 4 官方幻灯片。}
\end{figure}

\noindent\textbf{展开说明：}Set up a per-expert bias (making it more likely to get tokens) and use online learning

\begin{knowledgebox}{读图：Slide 42 应该怎么看}
这页是机制或证据页。先看它比较的是 compute、参数量、routing、负载均衡还是通信；再看结论支持的是质量提升、训练速度、推理成本还是系统复杂度。MoE 和 attention alternatives 的图常把模型结构和系统代价混在一起，读图时要分清“算法机制”和“硬件执行成本”。
\end{knowledgebox}


\begin{figure}[H]
\centering
\includegraphics[width=0.88\textwidth]{slide-042.jpg}
\caption{Slide 43: What happens when removing load balancing losses?}
\figsource{CS336 Spring 2026 Lecture 4 官方幻灯片。}
\end{figure}

\noindent\textbf{展开说明：}What happens when removing load balancing losses?

\begin{knowledgebox}{读图：Slide 43 应该怎么看}
这页是机制或证据页。先看它比较的是 compute、参数量、routing、负载均衡还是通信；再看结论支持的是质量提升、训练速度、推理成本还是系统复杂度。MoE 和 attention alternatives 的图常把模型结构和系统代价混在一起，读图时要分清“算法机制”和“硬件执行成本”。
\end{knowledgebox}

\subsection{本章小结}
本节的共同线索是 selective computation：通过结构选择让 token 不必总是访问所有历史或所有参数。但选择性越强，越需要额外机制保证信息流、负载均衡和训练稳定。

